\documentclass[12pt]{article}
\usepackage[utf8]{inputenc}
\usepackage[T1]{fontenc}
\usepackage[margin=1in]{geometry}
\usepackage{setspace}
\usepackage{xcolor}
\usepackage[colorlinks=true, urlcolor=blue!60!black, linkcolor=blue!60!black]{hyperref}
\setlength{\parindent}{2em}
\setlength{\parskip}{1.25em}
\renewcommand{\baselinestretch}{1.0}
\begin{document}
\begin{center}
{\LARGE \textit{The Civil Rights Cases}, 109 U.S. 3 (1883)\par}
\vspace{0.8em}
{\large Case Brief\par}
\end{center}
\vspace{1em}

\section{Memory Jogger}

In an 8--1 decision establishing the modern state action doctrine, the Supreme Court invalidated the public accommodations provisions of the Civil Rights Act of 1875, holding that the Fourteenth Amendment restricts only state action and that the Thirteenth Amendment does not reach private racial discrimination in inns, theaters, and railroads (\href{https://www.courtlistener.com/opinion/90897/civil-rights-cases/}{CourtListener}; \href{https://supreme.justia.com/cases/federal/us/109/3/}{Justia}).

\section{Detailed Case Facts}

On March 1, 1875, Congress enacted the Civil Rights Act of 1875, 18 Stat. 335, entitled ``An Act to protect all citizens in their civil and legal rights'' (\href{https://supreme.justia.com/cases/federal/us/109/3/}{Justia}). Section 1 declared that all persons within the jurisdiction of the United States were entitled to ``the full and equal enjoyment of the accommodations, advantages, facilities, and privileges of inns, public conveyances on land or water, theatres, and other places of public amusement,'' subject only to conditions established by law and applicable alike to citizens of every race and color, regardless of previous condition of servitude (\href{https://supreme.justia.com/cases/federal/us/109/3/}{Justia}). Section 2 made any violation by any person a federal misdemeanor punishable by a fine of \$500 to \$1,000 or imprisonment for thirty days to one year, and also created a \$500 civil penalty recoverable by the aggrieved party in an action of debt (\href{https://supreme.justia.com/cases/federal/us/109/3/}{Justia}).

Five separate disputes arising under Sections 1 and 2 were consolidated before the Supreme Court (\href{https://www.courtlistener.com/opinion/90897/civil-rights-cases/}{CourtListener}; \href{https://www.oyez.org/cases/1850-1900/109us3}{Oyez}):
\begin{itemize}
\item \textit{United States v. Stanley}: A federal indictment in Kansas charged Murray Stanley with denying a Black person the accommodations and privileges of an inn or hotel (\href{https://supreme.justia.com/cases/federal/us/109/3/}{Justia}).
\item \textit{United States v. Ryan}: A federal criminal information in California charged Michael Ryan with refusing a Black patron a seat in the dress circle of Maguire's Theatre in San Francisco (\href{https://supreme.justia.com/cases/federal/us/109/3/}{Justia}).
\item \textit{United States v. Nichols}: A federal indictment in Missouri charged Samuel Nichols with denying a Black person the accommodations and privileges of an inn or hotel (\href{https://supreme.justia.com/cases/federal/us/109/3/}{Justia}).
\item \textit{United States v. Singleton}: A federal indictment in New York charged Samuel Singleton with denying a patron (whose race was not specified in the indictment) full enjoyment of the Grand Opera House in New York City, allegedly not for reasons applicable alike to citizens of every race and color (\href{https://supreme.justia.com/cases/federal/us/109/3/}{Justia}).
\item \textit{Robinson \& Wife v. Memphis \& Charleston Railroad Co.}: Richard A. Robinson and his wife, Sallie Robinson (a Black citizen of Mississippi holding a ticket from Grand Junction, Tennessee, to Lynchburg, Virginia), brought a civil action in Tennessee for the \$500 statutory penalty after a railroad conductor excluded Mrs. Robinson from the ladies' car on account of her race (\href{https://supreme.justia.com/cases/federal/us/109/3/}{Justia}).
\end{itemize}

\section{Procedural History}

In \textit{Stanley} (Circuit Court for the District of Kansas), \textit{Nichols} (Circuit Court in Missouri), and \textit{Singleton} (Circuit Court for the Southern District of New York), the two circuit judges divided in opinion on the constitutionality of Sections 1 and 2 of the Civil Rights Act of 1875 and certified the question to the Supreme Court (\href{https://supreme.justia.com/cases/federal/us/109/3/}{Justia}; \href{https://www.govinfo.gov/content/pkg/USREPORTS-109/pdf/USREPORTS-109-3.pdf}{GovInfo}). In \textit{Ryan}, the Circuit Court for the District of California sustained a demurrer to the criminal information, and the United States sought review by writ of error (\href{https://supreme.justia.com/cases/federal/us/109/3/}{Justia}). In \textit{Robinson}, the Circuit Court for the Western District of Tennessee tried the case on the parties' shared assumption that the Act was valid; after the trial judge permitted evidence that the conductor suspected Mrs. Robinson of being an improper person because she was accompanied by a light-skinned young man he took to be white, the jury returned a verdict for the railroad, and the plaintiffs sought review by writ of error (\href{https://supreme.justia.com/cases/federal/us/109/3/}{Justia}).

The Solicitor General submitted \textit{Stanley}, \textit{Ryan}, \textit{Nichols}, and \textit{Singleton} together on November 7, 1882, with no appearances or briefs filed for the defendants; \textit{Robinson} was submitted on briefs on March 9, 1883 (\href{https://supreme.justia.com/cases/federal/us/109/3/}{Justia}). The Supreme Court decided all five consolidated cases on October 15, 1883 (\href{https://www.oyez.org/cases/1850-1900/109us3}{Oyez}; \href{https://supreme.justia.com/cases/federal/us/109/3/}{Justia}).

\section{Judicial Votes}

The Court decided the cases by an \textbf{8--1} vote (\href{https://www.oyez.org/cases/1850-1900/109us3}{Oyez}; \href{https://supreme.justia.com/cases/federal/us/109/3/}{Justia}):
\begin{itemize}
\item \textbf{Majority (8):} Justice Joseph P. Bradley (author), joined by Chief Justice Morrison R. Waite and Justices Samuel Freeman Miller, Stephen Johnson Field, William Burnham Woods, Stanley Matthews, Horace Gray, and Samuel Blatchford (\href{https://supreme.justia.com/cases/federal/us/109/3/}{Justia}; \href{https://www.oyez.org/cases/1850-1900/109us3}{Oyez}).
\item \textbf{Concurrences (0):} None.
\item \textbf{Dissent (1):} Justice John Marshall Harlan (\href{https://supreme.justia.com/cases/federal/us/109/3/}{Justia}; \href{https://www.oyez.org/cases/1850-1900/109us3}{Oyez}).
\end{itemize}

\section{Holding}

\textbf{Judgment:} Sections 1 and 2 of the Civil Rights Act of 1875 are unconstitutional and void as applied to the several States. In \textit{Ryan} and \textit{Robinson}, the judgments below were affirmed; in \textit{Stanley}, \textit{Nichols}, and \textit{Singleton}, the certified questions were answered that Sections 1 and 2 are unconstitutional and void, with judgment to be rendered on the indictments accordingly (\href{https://supreme.justia.com/cases/federal/us/109/3/}{Justia}).

\textbf{Rule Announced:} First, under the Fourteenth Amendment, only ``State action'' is prohibited; Section 5 empowers Congress solely to enact \textit{corrective} legislation counteracting or redressing unconstitutional state laws or state official acts, not primary and direct municipal legislation regulating private conduct (\href{https://supreme.justia.com/cases/federal/us/109/3/}{Justia}; \href{https://www.law.cornell.edu/constitution-conan/amendment-14/state-action-doctrine}{Constitution Annotated}). Second, although Section 2 of the Thirteenth Amendment authorizes primary and direct federal legislation against private individuals to abolish slavery and its ``badges and incidents,'' a private proprietor's denial of equal accommodations in inns, public conveyances, or places of public amusement does not impose a badge or incident of slavery (\href{https://supreme.justia.com/cases/federal/us/109/3/}{Justia}).

\textbf{Holding vs. Dicta:} The binding holding is confined to Congress's lack of authority under the Thirteenth and Fourteenth Amendments to enact Sections 1 and 2 of the 1875 Act as applied within the States. Justice Bradley explicitly reserved three questions as undecided: (1) whether equal access to public accommodations is a constitutional right demandable against state interference; (2) whether Sections 1 and 2 remained valid in the Territories or the District of Columbia, where Congress holds plenary municipal authority; and (3) whether Congress could regulate interstate public conveyances under the Commerce Clause (\href{https://supreme.justia.com/cases/federal/us/109/3/}{Justia}). Conversely, Bradley's observations that the Civil Rights Act of 1866 validly protected ``fundamental'' civil rights (contracting, suing, holding property) rather than ``social rights,'' and his remark that formerly enslaved persons must cease to be ``the special favorite of the laws,'' were broader dicta not strictly necessary to the state-action holding (\href{https://supreme.justia.com/cases/federal/us/109/3/}{Justia}).

\section{Analysis of Opinions}

\subsection{Majority Opinion (Justice Bradley)}

Justice Bradley began with the text of Section 1 of the Fourteenth Amendment (``No State shall...'') and reasoned that it prohibits ``State action of a particular character,'' not ``[i]ndividual invasion of individual rights'' (\href{https://supreme.justia.com/cases/federal/us/109/3/}{Justia}). Analogizing to the Contracts Clause of Article I, Section 10, Bradley argued that a constitutional prohibition against state wrongdoing empowers Congress under Section 5 only to enact \textit{corrective} remedies predicated on adverse state laws, customs having the force of law, or state official conduct (\href{https://supreme.justia.com/cases/federal/us/109/3/}{Justia}). Because Sections 1 and 2 of the 1875 Act applied directly to private individuals without reference to any state law or official failure---operating equally in states with fair public accommodations laws---the statute displaced state jurisprudence and created a federal ``code of municipal law,'' violating the Tenth Amendment (\href{https://supreme.justia.com/cases/federal/us/109/3/}{Justia}). Bradley contrasted these provisions with Section 4 of the same Act (upheld in \textit{Ex parte Virginia}, 100 U.S. 339 (1880), because jury exclusion was performed by state officers) and with the Civil Rights Act of 1866 as reenacted in 1870, whose penal provision targeted deprivations ``under color of'' state law or custom (\href{https://supreme.justia.com/cases/federal/us/109/3/}{Justia}).

Turning to the Thirteenth Amendment, Bradley conceded that it is self-executing and allows Congress to pass direct, primary legislation operating on private individuals to eradicate ``all badges and incidents of slavery'' (\href{https://supreme.justia.com/cases/federal/us/109/3/}{Justia}). However, he distinguished the ``inseparable incidents'' of slavery (compulsory service, inability to contract, sue, testify, or hold property) from ``social rights'' and private refusals of service in inns, theaters, and conveyances (\href{https://supreme.justia.com/cases/federal/us/109/3/}{Justia}). Treating every act of private racial discrimination as a badge of slavery, Bradley wrote, would be ``running the slavery argument into the ground'' (\href{https://supreme.justia.com/cases/federal/us/109/3/}{Justia}). Aggrieved citizens had to seek redress under state common law or state statutes unless state laws themselves denied equal protection (\href{https://supreme.justia.com/cases/federal/us/109/3/}{Justia}).

\subsection{Concurring Opinions}

No concurring opinions were filed (\href{https://supreme.justia.com/cases/federal/us/109/3/}{Justia}).

\subsection{Dissenting Opinion (Justice Harlan)}

Justice Harlan condemned the majority's analysis as ``entirely too narrow and artificial,'' arguing that it sacrificed the ``substance and spirit'' of the Reconstruction Amendments to ``subtle and ingenious verbal criticism'' (\href{https://supreme.justia.com/cases/federal/us/109/3/}{Justia}). Invoking \textit{Prigg v. Pennsylvania}, 41 U.S. (16 Pet.) 539 (1842), and the Fugitive Slave Acts of 1793 and 1850, Harlan pointed out a stark constitutional irony: before the Civil War, the Court had implied plenary, direct congressional power to regulate private individuals and override states to protect slaveholders' property rights, yet now denied Congress express enforcement power to protect freedom and citizenship (\href{https://supreme.justia.com/cases/federal/us/109/3/}{Justia}).

Under the Thirteenth Amendment, Harlan argued that because American slavery rested entirely on racial subordination, universal freedom necessarily entailed immunity from race-based discrimination in civil rights belonging to freemen (\href{https://supreme.justia.com/cases/federal/us/109/3/}{Justia}). He emphasized that railroads (which exercise eminent domain as public highways), common-law inns (which owe a public duty to receive all unobjectionable travelers), and licensed theaters (``clothed with a public interest'' under \textit{Munn v. Illinois}, 94 U.S. 113 (1877)) perform public or quasi-public functions and are not purely private enterprises (\href{https://supreme.justia.com/cases/federal/us/109/3/}{Justia}). Under the Fourteenth Amendment, Harlan argued that the Citizenship Clause of Section 1 is an \textit{affirmative} grant of national and state citizenship---not merely a prohibition on states---which Section 5 empowers Congress to protect through direct legislation (\href{https://supreme.justia.com/cases/federal/us/109/3/}{Justia}). Even on the majority's own state-action premise, Harlan maintained that railroads, innkeepers, and licensed amusement operators act as ``agents or instrumentalities of the State'' within \textit{Ex parte Virginia} (\href{https://supreme.justia.com/cases/federal/us/109/3/}{Justia}). Finally, he noted that \textit{Robinson} involved an interstate rail ticket from Tennessee to Virginia, squarely presenting the Commerce Clause question the majority declined to reach (\href{https://supreme.justia.com/cases/federal/us/109/3/}{Justia}).

\section{Examples: Future Applications}

\begin{enumerate}
\item \textbf{Same-Side Application (Private Violence Remedy Struck Down Under Section 5):} Congress enacts a federal civil remedy allowing victims of private, gender-motivated violent crimes to sue their attackers in federal court, citing Section 5 of the Fourteenth Amendment based on congressional findings of widespread state-court bias. \textit{Analysis:} Under \textit{The Civil Rights Cases}, the statute is unconstitutional under Section 5 because it is directed at private individuals rather than at state laws or state officials, and applies uniformly nationwide regardless of whether a particular state's courts have denied equal protection. This mirrors \href{https://supreme.justia.com/cases/federal/us/529/598/}{\textit{United States v. Morrison}, 529 U.S. 598 (2000)}.

\item \textbf{Same-Side Application (Purely Private Social Club Exclusion):} A privately owned, members-only dinner club that receives no state funding, holds no special public franchise, and operates on private property refuses to admit a prospective member solely on the basis of race, and the applicant sues the club directly under the Equal Protection Clause of the Fourteenth Amendment. \textit{Analysis:} Under \textit{The Civil Rights Cases}, the Fourteenth Amendment claim fails because ``individual invasion of individual rights'' unsupported by state law, judicial enforcement, or state authority is not state action.

\item \textbf{Opposite-Side Application (State Judicial Enforcement of Private Racial Covenants):} Private homeowners sign a restrictive covenant barring the sale of neighborhood homes to Black buyers; when a willing seller conveys a home to a Black family, neighbors obtain a state-court injunction enforcing the covenant. \textit{Analysis:} Although the private covenant standing alone does not violate the Fourteenth Amendment under \textit{The Civil Rights Cases}, the state court's affirmative judicial enforcement supplies the required state action and violates the Equal Protection Clause, as the Court later held in \href{https://supreme.justia.com/cases/federal/us/334/1/}{\textit{Shelley v. Kraemer}, 334 U.S. 1 (1948)}.

\item \textbf{Opposite-Side Application (Federal Public Accommodations Regulation Under the Commerce Clause or Thirteenth Amendment):} Congress enacts a statute barring racial discrimination by private motels serving interstate travelers (relying on the Commerce Clause) or prohibiting private racial discrimination in the sale and rental of housing (relying on Section 2 of the Thirteenth Amendment). \textit{Analysis:} Both statutes are valid without overruling the Fourteenth Amendment state-action holding of \textit{The Civil Rights Cases}. The Court expressly reserved Congress's Commerce Clause power in 1883 and later upheld Title II of the Civil Rights Act of 1964 in \href{https://supreme.justia.com/cases/federal/us/379/241/}{\textit{Heart of Atlanta Motel, Inc. v. United States}, 379 U.S. 241 (1964)}, and expanded Congress's power to define ``badges and incidents of slavery'' under the Thirteenth Amendment in \href{https://supreme.justia.com/cases/federal/us/392/409/}{\textit{Jones v. Alfred H. Mayer Co.}, 392 U.S. 409 (1968)}.

\item \textbf{Fence-Sitter (Privately Operated Venue Leased Inside a Municipal Building):} A private restaurant leases commercial space inside a city-owned parking garage constructed with public funds; the city maintains the building and benefits financially from the restaurant's patronage, while the restaurant refuses to serve Black customers without any city ordinance requiring segregation. \textit{Analysis:} This is a close case. Under a strict reading of Bradley's opinion, the restaurant owner is a private individual acting without a discriminatory state statute. Under a broader state-action or ``entwinement'' view echoing Harlan's dissent, the city's symbiotic financial relationship and public property ownership make the state a joint participant in the discrimination, blurring the boundary between private wrong and state sanction.
\end{enumerate}

\section{Critique}

\textbf{Originalist and Textualist Perspectives:} Defenders of \textit{The Civil Rights Cases} emphasize fidelity to the text of Section 1 of the Fourteenth Amendment (``No State shall...'') and structural federalism, arguing that reading Section 5 to authorize a general federal municipal code would obliterate the Tenth Amendment boundary between national and state authority (\href{https://constitution.heritage.org/essays/amdt-14-s1-b/}{Heritage Guide to the Constitution}). Libertarian-leaning scholars such as Richard Epstein have likewise noted that the state action requirement preserves a sphere of private property and freedom of association insulated from direct constitutional litigation (\href{https://scholarship.law.tamu.edu/cgi/viewcontent.cgi?article=1359&context=lawreview}{Texas A\&M Law Review}). Yet other originalist and historical scholars sharply challenge the decision's premise. Michael W. McConnell has demonstrated that the Reconstruction Congresses that drafted and immediately enforced the Fourteenth Amendment viewed its original meaning as far more robust than late-nineteenth-century judicial retrenchment suggested (\href{https://chicagounbound.uchicago.edu/cgi/viewcontent.cgi?article=12612&context=journal_articles}{McConnell, U. Chi. Law Unbound}). Similarly, Randy Barnett and Evan Bernick argue that the original meaning of ``equal protection of the laws'' obligated states to provide protection against private invasions of rights, meaning state \textit{inaction} or refusal to remedy private wrongs triggers federal enforcement power (\href{https://www.cato.org/blog/new-decision-considers-federal-power-criminalize-civil-rights-violations}{Cato Institute}). Historical-institutionalist Pamela Brandwein adds a nuanced revisionist reading: she contends that Bradley actually preserved a ``state neglect'' concept---allowing federal remedies when states fail to enforce neutral common-law duties---that later courts and lawyers mistakenly erased (\href{https://www.cambridge.org/core/journals/law-and-society-review/article/judicial-abandonment-of-blacks-rethinking-the-state-action-cases-of-the-waite-court/6EFE4577A05FAA2791AB1541BFBFC371}{Brandwein, Law \& Society Review}).

\textbf{Progressive and Doctrinal Perspectives:} Progressive constitutional scholars place \textit{The Civil Rights Cases} in the constitutional ``anti-canon'' alongside \textit{Plessy v. Ferguson} (\href{https://ideaexchange.uakron.edu/cgi/viewcontent.cgi?article=1127&context=conlawnow}{Akron ConLawNOW}). In his classic critique, Charles L. Black Jr. argued that the state action doctrine rests on an illusory public/private dichotomy because private property and exclusion are themselves backed by state trespass law and licensing (\href{https://scholarship.law.tamu.edu/cgi/viewcontent.cgi?article=1359&context=lawreview}{Texas A\&M Law Review}). Doctrinal critics also fault the Court for gutting the Thirteenth Amendment's promise of equal civil freedom and for forcing Congress eighty years later to rely on the Commerce Clause---an economic power---to address a fundamental moral and civil-rights wrong (\href{https://supreme.justia.com/cases/federal/us/379/241/}{Justia}).

\textbf{Underlying Competing Values:} The case pits \textit{federalism and private autonomy} (preserving state police power and a domain of private choice free from constitutional command) against \textit{substantive national equality and full citizenship} (ensuring that newly emancipated citizens are not relegated to second-class caste status by private monopolies, common carriers, and indifferent or hostile states).

\section{Key Quotations}

\begin{quote}
``It is State action of a particular character that is prohibited. Individual invasion of individual rights is not the subject matter of the amendment.''
\end{quote}
\noindent --- Justice Joseph P. Bradley, Opinion of the Court, 109 U.S. at 11 (\href{https://supreme.justia.com/cases/federal/us/109/3/}{Justia}; \href{https://www.courtlistener.com/opinion/90897/civil-rights-cases/}{CourtListener}).

\begin{quote}
``In fine, the legislation which Congress is authorized to adopt in this behalf is not general legislation upon the rights of the citizen, but corrective legislation, that is, such as may be necessary and proper for counteracting such laws as the States may adopt or enforce, and which, by the amendment, they are prohibited from making or enforcing, or such acts and proceedings as the States may commit or take, and which, by the amendment, they are prohibited from committing or taking.''
\end{quote}
\noindent --- Justice Joseph P. Bradley, Opinion of the Court, 109 U.S. at 13--14 (\href{https://supreme.justia.com/cases/federal/us/109/3/}{Justia}).

\begin{quote}
``When a man has emerged from slavery, and, by the aid of beneficent legislation, has shaken off the inseparable concomitants of that state, there must be some stage in the progress of his elevation when he takes the rank of a mere citizen and ceases to be the special favorite of the laws, and when his rights as a citizen or a man are to be protected in the ordinary modes by which other men's rights are protected.''
\end{quote}
\noindent --- Justice Joseph P. Bradley, Opinion of the Court, 109 U.S. at 25 (\href{https://supreme.justia.com/cases/federal/us/109/3/}{Justia}).

\begin{quote}
``The opinion in these cases proceeds, it seems to me, upon grounds entirely too narrow and artificial. I cannot resist the conclusion that the substance and spirit of the recent amendments of the Constitution have been sacrificed by a subtle and ingenious verbal criticism.''
\end{quote}
\noindent --- Justice John Marshall Harlan, Dissenting Opinion, 109 U.S. at 26 (\href{https://supreme.justia.com/cases/federal/us/109/3/}{Justia}).

\begin{quote}
``Today it is the colored race which is denied, by corporations and individuals wielding public authority, rights fundamental in their freedom and citizenship. At some future time, it may be that some other race will fall under the ban of race discrimination.''
\end{quote}
\noindent --- Justice John Marshall Harlan, Dissenting Opinion, 109 U.S. at 62 (\href{https://supreme.justia.com/cases/federal/us/109/3/}{Justia}).

\end{document}
